% main.tex — manuscript skeleton (venue-agnostic).
% Swap \documentclass for your venue style:
%   CVPR/ICCV : \documentclass[10pt,twocolumn,letterpaper]{article} + cvpr.sty
%   NeurIPS   : \documentclass{article} + neurips_2024.sty
%   ICLR      : iclr2025_conference.sty
%   Journal   : IEEEtran / elsarticle
% This default builds standalone with pdflatex + bibtex.
\documentclass[10pt,twocolumn,letterpaper]{article}

\usepackage[utf8]{inputenc}
\usepackage{graphicx}
\usepackage{booktabs}      % results tables (see latex-results-table skill)
\usepackage{amsmath,amssymb}
\usepackage{multirow}
\usepackage{xcolor}
\usepackage[colorlinks=true,allcolors=blue]{hyperref}
\usepackage{cleveref}

\graphicspath{{figures/}}

\title{[Paper Title]}
\author{[Authors]\\ [Affiliation]\\ \texttt{[email]}}
\date{}

\begin{document}
\maketitle

\begin{abstract}
Open-vocabulary object detectors (OVD) generalize poorly outside natural images, but
\emph{why} they fail in specialized domains is not measured: aggregate metrics such as
mAP collapse several distinct error sources into one number. We introduce
\textbf{OVD-Diagnose}, a diagnostic protocol and cross-domain benchmark that decomposes
OVD failure into three interpretable axes --- localization, semantic confusion, and
calibration --- using controlled inference modes (global vocabulary, oracle vocabulary,
and a class-agnostic proposer). Applied to frozen detectors across specialized domains,
the protocol reveals that the dominant failure axis is \emph{domain-specific}: on aerial
imagery models localize objects but are bottlenecked by vocabulary confusion (oracle
prompting more than triples AP), whereas on chest radiographs localization itself fails
(a class-agnostic SAM proposer recovers under $1\%$ of findings), rendering vocabulary
moot. We further show that semantic confusion grows with detector capability. These
patterns are invisible to closed-set error tools, and the protocol is training-free,
model-agnostic, and reproducible on modest hardware.
\end{abstract}

\section{Introduction}
\label{sec:intro}
Open-vocabulary object detection (OVD)~\cite{radford2021clip} promises detection of
arbitrary categories from text, yet its zero-shot performance degrades sharply outside
the natural-image distribution it was trained on. This matters precisely where OVD would
be most useful: aerial survey, medical screening, and other specialized domains where
exhaustively annotating every category is impractical. Recent benchmarks confirm the
degradation but report it as a single collapsed number (mAP or F1), which cannot say
\emph{why} a detector fails --- whether it cannot find the objects, cannot name them, or
simply mis-thresholds its own confidence. These are different problems demanding different
remedies, and conflating them has slowed progress on specialized-domain OVD.

We argue that OVD failure should be \emph{decomposed}, not summarized. Closed-set error
toolkits (e.g.\ TIDE) dissect detector errors but predate and ignore the axis unique to
OVD: error induced by the open vocabulary itself. We introduce \textbf{OVD-Diagnose}, a
training-free protocol that isolates three axes with controlled inference modes ---
prompting a detector with the full domain vocabulary, with only the classes present in an
image, and (for localization) with a class-agnostic proposer --- and reads off
localization ($L$), semantic-confusion ($S$), and calibration ($C$) error from the gaps
between them.

Our contributions are:
\begin{itemize}
  \item \textbf{A diagnostic protocol} that decomposes OVD failure into localization,
  semantic-confusion, and calibration axes, validated by synthetic controls and an
  IoA-F1 anchor that reproduces prior benchmarks' model ordering.
  \item \textbf{A cross-domain benchmark} over frozen state-of-the-art OVD detectors and
  specialized domains, released as an open toolkit for arbitrary detector/domain pairs.
  \item \textbf{Findings:} the dominant failure axis is domain-specific (aerial is
  vocabulary-bound; medical is localization-bound), and semantic confusion increases with
  detector capability --- neither observable from aggregate mAP.
\end{itemize}

\section{Related Work}
\label{sec:related}
[Group by theme: detection~\cite{ren2015faster,carion2020detr,zhang2023dino},
backbones~\cite{he2016resnet,dosovitskiy2021vit,liu2021swin},
multimodal~\cite{radford2021clip,li2023blip2,liu2023llava}. Position your work.]

\section{Method}
\label{sec:method}
[Formal setup. Input $x \in \mathbb{R}^{H\times W\times 3}$, prediction $\hat{y}$,
ground truth $y$. Describe architecture, then loss.]
\begin{equation}
  \mathcal{L} = \mathcal{L}_{\text{cls}} + \lambda \mathcal{L}_{\text{reg}}.
  \label{eq:loss}
\end{equation}
[Reference figure with \cref{fig:arch}.]
% \begin{figure}[t]\centering
%   \includegraphics[width=\linewidth]{arch}
%   \caption{Architecture overview.}\label{fig:arch}
% \end{figure}

\section{Experiments}
\label{sec:exp}
\subsection{Setup}
[Datasets, splits, metrics (exact names), implementation details, hardware.]

\subsection{Main Results}
% Cross-domain OVD-Diagnose fingerprint — generated from runs/diag/*/fingerprints.csv
% Regenerate: python tools/make_paper_table.py --domains aerial:... medical:... --out paper/tables/ovd_fingerprint.tex
% Aerial: 1325 imgs (partial split). Medical (VinDr): 200 imgs. Update after full runs.
\begin{table}[t]
  \centering
  \caption{Cross-domain OVD failure fingerprint. AP is COCO mAP@[.5:.95].
  $L{=}1{-}\mathrm{AR}_{\mathrm{SAM}}$ (domain-level localizability).
  $S_{\mathrm{norm}}$ is the fraction of achievable AP lost to vocabulary confusion;
  $^{\dagger}$ unstable when AP$_{\mathrm{oracle}}{\approx}0$.
  $C_{\mathrm{ece}}$ is expected calibration error. Aerial is \emph{semantic-confusion
  bound} (models localize; oracle triples AP); medical is \emph{localization bound}
  ($L{\approx}1$: even SAM boxes $<$1\% of findings).}
  \label{tab:fingerprint}
  \begin{tabular}{ll cc c c c c}
    \toprule
    Domain & Model & AP$_{\mathrm{g}}$ & AP$_{\mathrm{o}}$ & $L$ & $S_{\mathrm{norm}}$ & $C_{\mathrm{ece}}$ & IoA-F1 \\
    \midrule
    \multirow{3}{*}{Aerial} & YOLO-World    & 0.004 & 0.018 & 0.86 & 0.76 & 0.009 & 0.003 \\
                            & OWLv2         & 0.015 & 0.079 & 0.86 & 0.81 & 0.024 & 0.039 \\
                            & Grounding DINO & 0.004 & 0.050 & 0.86 & 0.92 & 0.084 & 0.006 \\
    \midrule
    \multirow{3}{*}{Medical} & YOLO-World    & 0.000 & 0.000 & 1.00 & 0.91$^{\dagger}$ & 0.001 & 0.000 \\
                             & OWLv2         & 0.002 & 0.004 & 1.00 & 0.48$^{\dagger}$ & 0.011 & 0.072 \\
                             & Grounding DINO & 0.000 & 0.001 & 1.00 & 0.95$^{\dagger}$ & 0.085 & 0.082 \\
    \bottomrule
  \end{tabular}
\end{table}

\Cref{tab:fingerprint} reports the failure fingerprint of five frozen open-vocabulary
detectors across two specialized domains under our protocol.

\subsection{Cross-domain findings}
\label{sec:findings}
Two observations dominate. \textbf{(1) The bottleneck axis is domain-specific.} On
aerial imagery, models can localize (SAM recovers $14\%$ of objects, and prompting with
the oracle vocabulary triples AP), so failure is driven by \emph{semantic confusion}:
$S_{\mathrm{norm}}$ reaches $0.76$--$0.92$. On chest radiographs the picture inverts:
$L{\approx}1.0$ --- even a class-agnostic SAM proposer boxes under $1\%$ of findings ---
and providing oracle classes does not help, because \emph{localization}, not vocabulary,
is the wall. Thus aerial is semantic-confusion bound while medical is localization bound,
despite an identical protocol. \textbf{(2) Semantic confusion scales with model
capability.} Within the aerial domain, $S_{\mathrm{norm}}$ rises monotonically from
YOLO-World ($0.76$) to OWLv2 ($0.81$) to Grounding DINO ($0.92$): stronger detectors
extract usable features but are increasingly swamped by the open vocabulary. These
patterns are invisible to closed-set error tools (e.g.\ TIDE), which cannot isolate the
vocabulary-induced axis.

\subsection{Validation}
\label{sec:validation}
The IoA-F1 anchor reproduces the model ordering of prior aerial benchmarks (OWLv2 best),
confirming the harness. Synthetic controls (Sec.~\ref{sec:method}) verify that each axis
responds only to its intended perturbation.

\section{Conclusion}
\label{sec:conc}
[Summary, limitations, future work.]

{\small
\bibliographystyle{ieee_fullname}   % or plainnat / abbrvnat per venue
\bibliography{refs}
}

\end{document}
